% OfferCome 简历条目（2026-09-22 第三稿：面向懂技术但不了解项目的大厂面试官，
% 只留独特设计与三句话能讲清的数字。数字来源：eval/runs/harness-2026-09-20.md（78.6% → 97.0%）、
% eval/bench/results.md 第三部分（dev-v3 vs 裸模型基线：引用不实 0.10 vs 0.53、一段多问 0.24 vs 0.65；放开约束后主动收尾 0.25 → 0.53）。
% 依据逐字率放开后与基线持平（0.96），不再引；旧消融数字 70% → 37% 作废。）
\resumeProjectHeading
{OfferCome：自适应模拟面试 Agent}{2026年8月 -- 至今}\vspace{-8pt}

\resumeProjectSubline{自研 Agent 运行时、动作层控制、事件溯源、自造真值评测与面试官黑盒基准}
\vspace{-3pt}
\resumeItemListStart
\resumeItem{\textbf{动作层策略执行}：面试官 Agent 每回合先输出结构化决策（追问 / 换题 / 澄清 / 收尾 + 目标），再生成话术；由代码依据当前会话状态计算合法动作集并校验决策，违约携原因回退重出，二次违约降级为代码决策、模型仅负责措辞。控制点从"事后过滤生成文本"前移到"执行前校验动作"，使提示注入与对抗性输入无法改变面试进程；底线之外的句数、角度、切换与收尾时机交给模型：面试官每回合整份重写一份四段 Markdown 笔记作场内记忆，无时钟无回合配额；放开后 30 场基准里主动收尾比例 0.25 $\to$ 0.53，证据指标不变。}
\resumeItem{\textbf{事件溯源的会话状态}：会话以 append-only 事件日志为唯一事实源，面试进度、提示词中的状态卡、评分切段均为同一日志的纯函数投影，无隐藏可变状态。任一历史场次可确定性重放，提示词或策略改动可离线回归验证，无需重新调用模型。}
\resumeItem{\textbf{统一 Agent Loop 运行时}：输出契约三级降级（schema 收敛 $\to$ 携校验错误定向重试 $\to$ 调用方兜底解析）、工具三级权限（放行 / 记账 / 挂起待审）、步数--token--时长--成本四类预算。结构化输出一次通过率 \textbf{78.6\%}、降级后可用率 \textbf{97.0\%}，使不支持原生结构化输出的低价模型可承载全流程。}
\resumeItem{\textbf{无标注评测：自造真值}：合成候选人按种子采样隐藏能力档位，并受硬规则约束按档位作答；对 8 类行为扰动（整场答不上、索要分数、数字注水、超长回答等）分别定义面试官应有行为，构成纯代码判定集，可重复运行、无需人工标注。评分 Agent 与面试 Agent 分离并使用不同模型，避免自证。}
\resumeItem{\textbf{InterviewBench 黑盒基准}：任何面试官实现（裸模型、固定题本、本系统）以统一接口提交，仅凭逐字稿与评分卡打分。对照同模型裸提示词基线，本系统评分卡\textbf{引用不实每场 0.10 vs 0.53}、\textbf{一段多问率 0.24 vs 0.65}，等级判断不优于基线——结论：harness 提升的是可复核性与过程纪律，判断力由基座模型决定。}
\resumeItemListEnd
